%% Template: modern-navy CV
%% Comprehensive master reference, profile-agnostic ([PLACEHOLDER] tokens).
%% German-language design: clean sans-serif, navy section rules, two-column
%% skills table. Standalone article-class design -- NOT moderncv.
%%
%% Compile with: lualatex template.tex

\documentclass[10pt,a4paper]{article}
\usepackage{fontspec}
\setmainfont{TeX Gyre Heros}
\usepackage[a4paper,margin=1.15cm,top=0.9cm,bottom=0.9cm]{geometry}
\usepackage{xcolor}
\definecolor{navy}{HTML}{1B3A5C}
\definecolor{gray}{HTML}{555555}
\usepackage{hyperref}
\hypersetup{colorlinks=true,linkcolor=navy,urlcolor=navy,pdftitle={[YOUR_NAME] - Lebenslauf}}
\usepackage{array}

\setlength{\parindent}{0pt}
\setlength{\parskip}{0pt}
\pagenumbering{gobble}

% German text has long unhyphenated compounds and no German hyphenation
% patterns are loaded (avoiding babel/polyglossia - both trigger a MiKTeX
% on-the-fly install hang for the German .ldf/pattern files in this
% sandbox); allow extra stretch instead so lines don't overflow the margin.
\tolerance=2500
\emergencystretch=2.5em
\hbadness=10000

% Compact list spacing (plain article \@listI -- no enumitem dependency)
\makeatletter
\renewcommand\@listI{\leftmargin\leftmargini
  \topsep=1pt \itemsep=0pt \parsep=0pt}
\@listi
\makeatother

% ---- section formatting: bold navy label, own line, full-width rule below ----
\newcommand{\cvsection}[1]{%
  \par\vspace{3pt}
  {\bfseries\color{navy}\fontsize{10.5}{12}\selectfont #1}\par
  \vspace{0pt}
  \noindent{\color{navy}\rule{\linewidth}{0.9pt}}\par
  \vspace{2pt}}

% ---- entry helper: title/date on one line, subtitle/location on the next ----
\newcommand{\entryhead}[2]{%
  {\bfseries #1}\hfill{\color{gray}#2}\par}
\newcommand{\entrysub}[2]{%
  {\itshape\color{gray}\small #1}\hfill{\itshape\color{gray}\small #2}\par\vspace{1pt}}

\begin{document}

% ============================================================
%     HEADER
% ============================================================
{\fontsize{25}{27}\selectfont\bfseries\color{navy}[YOUR_NAME]}\par
\vspace{2pt}
{\color{gray}\large [Your Field] \textbar\ [Your Specializations]}\par
\vspace{4pt}
{\small [Your Address, City, Postcode] \textbar\ [+XX XXXXXXXXXX] \textbar\ \href{mailto:[your.email@example.com]}{[your.email@example.com]}\\
\href{[https://linkedin.com/in/your-profile]}{[linkedin.com/in/your-profile]} \textbar\ [Optional: date/place of birth, nationality, or other line German CVs conventionally include]}\par
\vspace{2pt}
\noindent{\color{navy}\rule{\linewidth}{1.1pt}}

% ============================================================
%     PROFIL
% ============================================================
\cvsection{Profil}
\small
[Write a 3-5 sentence profile statement tailored to the role, in German. Focus on what makes you uniquely qualified: field of study, hands-on experience, and the kind of ownership you like to take.]

% ============================================================
%     BERUFSERFAHRUNG
% ============================================================
\cvsection{Berufserfahrung}
\small
\entryhead{[Employer] \textperiodcentered\ [Job Title]}{[MM/YYYY -- MM/YYYY]}
\entrysub{[Team/Department]}{[City]}
\begin{itemize}
  \item [Achievement/responsibility bullet 1 -- be concrete, quantify where possible]
  \item [Achievement/responsibility bullet 2]
  \item [Achievement/responsibility bullet 3]
\end{itemize}
\vspace{2pt}
\textbf{[Sub-project or notable initiative title]} -- [one-line description]
\begin{itemize}
  \item [Detail bullet 1]
  \item [Detail bullet 2]
\end{itemize}

% ============================================================
%     PROJEKTE
% ============================================================
\cvsection{Projekte}
\small
\entryhead{[Project Name] (\href{[https://github.com/your-username/project]}{GitHub})}{[MM/YYYY]}
\entrysub{[Personal Project \textperiodcentered\ Tech stack]}{}
\begin{itemize}
  \item [What you built and why]
  \item [A concrete, quantified result]
\end{itemize}
\vspace{3pt}

\entryhead{[Project Name 2]}{[MM/YYYY -- MM/YYYY]}
\entrysub{[Institution/Context]}{}
\begin{itemize}
  \item [Bullet 1]
  \item [Bullet 2]
\end{itemize}

% ============================================================
%     AUSBILDUNG
% ============================================================
\cvsection{Ausbildung}
\small
\entryhead{[Institution]}{[MM/YYYY -- MM/YYYY]}
\entrysub{[Degree, Field]}{}
\begin{itemize}
  \item [Grade / honors / scholarship, if relevant]
\end{itemize}

% ============================================================
%     AUSSERSCHULISCHES ENGAGEMENT
% ============================================================
\cvsection{Ausserschulisches Engagement}
\small
\entryhead{[Organization] \textperiodcentered\ [Role]}{[MM/YYYY -- MM/YYYY]}
\entrysub{[Context]}{[City]}
\begin{itemize}
  \item [Bullet 1]
\end{itemize}

% ============================================================
%     TECHNISCHE FÄHIGKEITEN
% ============================================================
\cvsection{Technische Fähigkeiten}
\small
\begin{tabular}{@{}>{\bfseries}p{4.4cm}p{12.9cm}@{}}
[Category 1] & [Skills, tools, frameworks] \\[2pt]
[Category 2] & [Skills, tools, frameworks] \\[2pt]
[Category 3] & [Skills, tools, frameworks] \\
\end{tabular}

% ============================================================
%     SPRACHEN
% ============================================================
\cvsection{Sprachen}
\small
\textbf{[Language 1]} [Level] \textbar\ \textbf{[Language 2]} [Level] \textbar\ \textbf{[Language 3]} [Level]

\end{document}
